\begin{table}[!h]
\centering
\caption{\label{tab:tab:race_census_comparison}Racial and Ethnic Composition of the Survey Sample and 2020 U.S. Population}
\centering
\begin{threeparttable}
\begin{tabular}[t]{lrrrr}
\toprule
Race/Ethnicity & Sample $N$ & Sample \% & 2020 U.S. Population (millions) & 2020 Census \%\\
\midrule
Asian & 78 & 4.9\% & 24.0 & 7.2\%\\
Black or African American & 266 & 16.6\% & 46.9 & 14.2\%\\
Hispanic or Latino & 237 & 14.8\% & 62.1 & 18.7\%\\
Middle Eastern or North African & 6 & 0.4\% & 3.5 & 1.1\%\\
American Indian or Alaska Native & 58 & 3.6\% & 9.7 & 2.9\%\\
\addlinespace
Native Hawaiian or Other Pacific Islander & 5 & 0.3\% & 1.6 & 0.5\%\\
White & 1179 & 73.7\% & 235.4 & 71.0\%\\
\bottomrule
\end{tabular}
\begin{tablenotes}
\item \textit{Note:} 
\item Sample racial categories are non-mutually exclusive; respondents could identify with more than one race. Accordingly, percentages need not sum to 100 percent. Census race estimates report each racial group alone or in combination with one or more other races. Hispanic or Latino origin is measured by the Census as an ethnicity and includes persons of any race. In the 2020 Census, Middle Eastern or North African (MENA) was not a separate race category and was formally classified within White; the MENA estimate shown here uses detailed Census responses identifying MENA origin and therefore overlaps with the Census White category. Census benchmarks refer to the total 2020 U.S. population.
\end{tablenotes}
\end{threeparttable}
\end{table}
